\chapterhead{50}{ORR-11}{Case Study: Financial Crime and Fraud}{2}{2}

\lohead{50}{a}{Describe elements of a control framework to manage financial fraud
risk and money laundering risk.}

Financial crime encompasses internal and external fraud, money laundering (ML),
and terrorism financing (TF), with definitions provided by authoritative bodies
like the FCA and BCBS. Effective management of these risks involves comprehensive
frameworks for internal and external fraud management, as well as anti-money
laundering (AML) measures.

\begin{orkeybox}
{\sffamily\bfseries\fontsize{8.6}{10}\selectfont\color{navy}Internal versus
external fraud}\par
\vspace{2pt}
Internal fraud involves dishonest acts by employees, including unauthorised
activities and theft, while external fraud is committed by third parties,
including theft, forgery, and cybersecurity breaches.\par
\vspace{7pt}
{\sffamily\bfseries\fontsize{8.6}{10}\selectfont\color{navy}Money laundering (ML)
and terrorism financing (TF)}\par
\vspace{2pt}
Money laundering is the process of concealing the origins of illegally obtained
money; terrorism financing refers to the provision of funds for terrorist
activities.
\end{orkeybox}

\orsub{Financial crime risk management}

\orsubb{1) Internal controls}
\begin{lettered}
\item \textbf{Selection:} proper screening of employees and third parties during
onboarding.
\item \textbf{Prevention:} segregating duties, access controls, and clear
authorisation processes.
\item \textbf{Detection:} strong supervision, monitoring, reconciliation, and
reporting of unusual activity.
\item \textbf{Deterrence:} sanctions, lawsuits, and disciplinary actions against
offenders.
\end{lettered}

\orsubb{2) External fraud management}

Similar to internal fraud management, but focused on external threats like
customers, business partners, and criminals. Techniques include security
personnel, partnering with law enforcement, and specialised teams to combat
fraud.

\orsubb{3) AML risk management}

It involves four control steps:

\begin{lettered}
\item \textbf{Selection:} select customers through KYC, due diligence, verifying
identity and fund origins.
\item \textbf{Prevention:} implement risk-based controls with employee training.
Classify customers by occupation, activity, and so on.
\item \textbf{Detection:} strong governance, transaction monitoring systems, and
alert systems for abnormal activity.
\item \textbf{Deterrence:} lawsuits, account closure, and reporting suspicious
activity to authorities deter money laundering.
\end{lettered}

\orsub{The three phases of money laundering}

\begin{center}
\begin{tikzpicture}[font=\fontsize{8.2}{10.4}\selectfont,
  st/.style={text width=4.2cm,align=center,rounded corners=3pt,inner sep=6pt,
             text=white,font=\sffamily\bfseries\fontsize{8.8}{10.6}\selectfont},
  sb/.style={text width=4.2cm,align=left,rounded corners=3pt,inner sep=7pt,
             fill=paleblue,draw=palenavy,line width=0.6pt,text=navy}]
\foreach \i/\x/\c/\t/\d in {%
  1/-5.6/navy/{1. Placement}/{Criminals disguise illegal funds' origin by placing them in the financial system (cash transfers, false invoicing), often by \textbf{smurfing}: a series of small amounts into the banking system.},
  2/0/teal/{2. Layering}/{Criminals obscure the source of funds through complex transactions across multiple accounts.},
  3/5.6/forest/{3. Integration}/{Criminals integrate laundered funds back into the economy (asset purchases, fake payments).}}
{
  \node[st,fill=\c,anchor=north] (t\i) at (\x,0) {\t};
  \node[sb,anchor=north] (b\i) at (t\i.south) {\d};
}
\foreach \a/\b in {1/2,2/3}
  \draw[-{Stealth[length=5pt]},navy!35,line width=1.4pt]
    ($(t\a.east)+(0.05,0)$) -- ($(t\b.west)+(-0.05,0)$);
\end{tikzpicture}
\end{center}
\orfigcap{Placement, layering and integration.}

\begin{ornotebox}{Also worth knowing}
New RegTech companies with automation and machine learning are speeding up Know
Your Customer (KYC) processes, but regulators are checking their effectiveness.
Financial institutions must also screen clients against sanctions lists like those
from the US Treasury Department's Office of Foreign Assets Control (OFAC).
\end{ornotebox}

% ---------------------------------------------------------------------
\lohead{50}{b}{Summarize the regulatory findings and describe the lessons learned
from the USAA case study.}

\orsub{USAA Federal Savings Bank case study}

\begin{checks}
\item USAA Federal Savings Bank had a weak AML program for at least five years
(2016--2021).
\item The program failed to meet minimum requirements of the Bank Secrecy Act
(BSA) and Anti-Money Laundering (AML) standards.
\item The bank inadequately monitored suspicious activities and failed to report
them to regulators.
\end{checks}

\begin{orexambox}{Lessons learned}
\begin{itemize}
\item Financial institutions need strong AML controls and well-designed risk
management systems.
\item Staff assigned to AML compliance must be properly trained.
\item Automated monitoring systems should be effective and not overwhelm with
false positives.
\item The number of alerts generated by such systems needs to be manageable.
\item Increased remote transactions due to the pandemic highlight the importance
of adequate AML controls.
\end{itemize}
\end{orexambox}
